\subsection{Baptism: Regeneration into the Paschal Body}

\begin{studybox}{Governing thesis}
Baptism is the sacrament of faith and new birth in which Christ, through true water and the Trinitarian word, washes a human person into his death and resurrection. It remits every sin and every punishment due to sin, gives sanctifying grace and the infused virtues, incorporates the recipient into the Church, and imprints the baptismal character. It is the necessary sacramental gate to Christian worship, ordered from its first moment toward Confirmation, the Eucharist, and the unveiled life of the Trinity.
\end{studybox}

\subsubsection{Water before the font: creation, judgment, passage, and promise}

The baptismal water belongs first to God, not to religious invention. At creation the Spirit of God hovers over the waters before light and ordered life emerge (Gen. 1:2). Water is therefore already capable of bearing two related meanings: the unmastered deep from which no creature can rescue itself, and the fruitful element over which the divine breath brings an ordered world. Baptism does not declare matter evil and escape it. The Creator takes a creature he made, joins it to his word, and makes bodily washing the instrument of supernatural generation.

The flood adds judgment and new beginning. A violent world passes under water; Noah's household is carried through to a renewed earth under covenant and the sign of peace. Peter does not reduce the type to moral cleansing: the ark's deliverance corresponds to Baptism, ``not as a removal of dirt from the body but as an appeal to God for a clear conscience, through the resurrection of Jesus Christ'' (1~Pet. 3:20--21). The contrast guards sacramental realism. Baptism truly uses bodily water, yet its proper effect exceeds every physical power because the risen Christ acts through it.

At the Red Sea, Israel is freed from bondage by passing through waters that overwhelm the oppressor. Paul calls this a baptism ``into Moses'' in cloud and sea, then warns that sacramental privilege without persevering fidelity did not save the rebellious (Exod. 14; 1~Cor. 10:1--12). The type thus joins objective gift and moral summons. The font truly transfers dominion---from Adam's solidarity to Christ's---but the liberated person must walk toward the promised inheritance rather than turn desire back toward Egypt.

\enlargethispage{\baselineskip}
Jordan continues the pattern. Israel crosses it into the land; Naaman descends seven times and rises cleansed; Elisha's iron rises from its depths; Ezekiel sees purifying water promised with a new heart and Spirit, and later a river flowing from the temple to make the dead sea live (Josh. 3; 2~Kgs. 5; 6:1--7; Ezek. 36:24--28; 47:1--12). These events do not independently predict every baptismal detail. Together they establish a providential grammar: descent and rising, cleansing and adoption, judgment and mercy, Spirit and water, temple and life.

Circumcision supplies another covenant line. It marks incorporation, promise, and a people before the recipient can choose his ancestry. Paul joins a circumcision ``made without hands'' to burial and rising with Christ in Baptism (Col. 2:11--13). Baptism is not merely the Christian replacement of one ethnic boundary marker; it fulfills covenant incorporation by joining Jew and Gentile, slave and free, male and female to Christ (Gal. 3:26--29). Yet the analogy helps explain why divine gift can precede reflective choice in an infant: grace creates belonging and summons the freedom it heals.

\begin{threestable}{Figure}{Baptismal fulfillment}{Necessary guardrail}
Creation's waters and Spirit & The same Creator uses water and word for a new creation in the Spirit. & The element has no autonomous magical power; its sacramental instrumentality comes from Christ's institution.\\
Flood and ark & Sin is judged, a household is borne through death, and a covenant world begins. & Peter distinguishes sacramental appeal and resurrection from ordinary bodily washing.\\
Sea and Jordan & Bondage is broken by passage; the people enters covenant pilgrimage and inheritance. & Initiation is not final perseverance apart from faith, hope, charity, and obedience.\\
Circumcision and prophetic washing & God marks a people, promises a new heart, and joins water to the Spirit. & Baptism universalizes adoption in Christ; it does not reproduce ethnic exclusion.\\
\end{threestable}

\subsubsection{Christ enters the waters and opens his side}

The sinless Son receives John's baptism of repentance not because he needs cleansing, but to ``fulfill all righteousness'' and identify himself with the people he will redeem (Matt. 3:13--17). The opened heavens, paternal voice, and descending Spirit disclose a Trinitarian pattern that the final command will make explicit. The Fathers often say that Christ sanctifies the waters by entering them. The statement is not a claim that every river became a sacrament independently of ecclesial form; it confesses that the incarnate Word makes creaturely water apt to serve his saving humanity.

Jesus tells Nicodemus that one must be born ``of water and Spirit'' to enter God's kingdom (John 3:3--5). The dialogue holds together gift from above and sensible means below. Natural generation cannot produce divine filiation, and human effort cannot command the Spirit; nevertheless the promised birth is not an invisible conversion opposed to water. John's Gospel proceeds through living water, the healing pool, the washing of the disciples' feet, and finally the blood and water flowing from Christ's pierced side (John 4; 5; 7:37--39; 13; 19:34). The side of the sleeping new Adam gives the Fathers an image of the Church born in Baptism and Eucharist.

After the Resurrection the Lord sends the Eleven to make disciples of all nations, ``baptizing them in the name of the Father and of the Son and of the Holy Spirit,'' and teaching obedience to all he commanded (Matt. 28:18--20). The command supplies sacramental universality, Trinitarian form, ecclesial mission, and lifelong catechesis. Mark's longer ending joins Baptism and faith; Luke and Acts join repentance, forgiveness, the Spirit, and entry into the community. No single text should be isolated from the Paschal whole. Christ institutes Baptism through his teaching, death, resurrection, and commission; Pentecost makes the apostolic Church its ministerial body.

\subsubsection{Apostolic proclamation: death, clothing, body, and bath}

At Pentecost Peter answers convicted hearers: ``Repent, and be baptized every one of you in the name of Jesus Christ for the forgiveness of your sins; and you shall receive the gift of the Holy Spirit'' (Acts 2:38). ``In the name of Jesus'' identifies allegiance, authority, and the saving Lord; it need not be read as a surviving non-Trinitarian liturgical formula set against Matthew 28. Acts repeatedly joins preaching, faith, water, household incorporation, and ecclesial communion: Samaritans, the Ethiopian, Saul, Cornelius, Lydia's household, the Philippian jailer's household, and the Ephesians enter through a visible apostolic act (Acts 8--10; 16; 19).

Romans gives the controlling Paschal ontology. Those baptized into Christ Jesus are baptized into his death, buried with him, and raised so that they may walk in newness of life (Rom. 6:3--11). The rite does not stage an absent drama as devotional theater. Christ's once-for-all death remains historically unrepeatable, while its saving power reaches this person sacramentally. The baptized is transferred into a new solidarity: dead to sin's dominion, alive to God, still obliged to yield bodily members as instruments of righteousness.

Paul adds complementary images. Baptism is clothing with Christ and abolition of covenant inferiority without erasing creaturely embodiment (Gal. 3:27--29). In one Spirit all are baptized into one Body (1~Cor. 12:13). Christ cleanses the Church ``by the washing of water with the word'' to present her holy and without blemish (Eph. 5:25--27). Titus names a washing of regeneration and renewal in the Holy Spirit, poured out through Jesus Christ so that the justified become heirs in hope (Titus 3:4--7). Hebrews speaks of enlightened Christians whose bodies are washed with pure water; First John joins water, blood, and Spirit in testimony (Heb. 10:22; 1~John 5:6--8).

These images cannot be collapsed into a single metaphor precisely because Baptism causes a rich reality. It is bath, birth, burial, resurrection, clothing, enlightenment, sealing, incorporation, nuptial cleansing, adoption, and Exodus. The metaphors converge on a real change of relation and principle of life; none authorizes the claim that the external rite is sufficient while its ecclesial and moral form is denied.

\subsubsection{Eastern and Western Fathers at the one font}

The \emph{Didache} gives the earliest extra-biblical ritual directions: baptize in the triune name, preferably in living water, otherwise in other water, and if sufficient water is lacking pour water three times upon the head (7.1--4). The hierarchy of modes shows both reverence for immersion's expressive fullness and a sober recognition that sacramental identity does not depend on abundance. It is strong evidence for Trinitarian form and for pouring, not a complete theory of validity.

Justin Martyr describes converts brought to water and regenerated ``in the same manner in which we were ourselves regenerated,'' invoking Father, Son, and Holy Spirit (\emph{First Apology} 61). He connects the bath to deliberate faith, repentance, illumination, and Isaiah's call to washing. His adult missionary setting should not be made into a denial of infant Baptism; it shows how the Church initiates catechumens capable of hearing and choosing.

Irenaeus places regeneration within recapitulation: the Word who passed through every age sanctifies human life and gives the Spirit who renews the creature (\emph{Against Heresies} III.17; V.15). Tertullian marvels that water, a simple and abundant element, becomes the vehicle of so great a work after invocation of God, while insisting that the flesh is washed so the soul may be cleansed (\emph{De baptismo} 1, 4, 7--8). His later preference to delay certain baptisms is a disciplinary judgment, not denial of baptismal regeneration.

Cyril of Jerusalem addresses the newly baptized as persons who went down into the water and in a figure were buried with Christ, emerging into newness of life (\emph{Mystagogical Catecheses} II.4--7). He refuses both bare symbolism and crude physicalism: the bodily action and saving reality are distinct yet united by divine promise. Gregory Nazianzen calls Baptism gift, grace, anointing, enlightenment, garment of incorruption, bath of regeneration, and seal (\emph{Oration} 40.3--4). His rhetoric registers the one effect under many scriptural aspects.

Ambrose reads the Spirit over creation's waters, the flood, Red Sea, Marah, Naaman, and Christ's Baptism as a pedagogy toward the Milanese font (\emph{De mysteriis} 3--4). Augustine's anti-Donatist controversy supplies a decisive causal clarification. Sacraments belong to Christ even when an unworthy minister handles them; the minister's sin may condemn the minister but cannot make Christ's Baptism private property. ``When Peter baptizes, Christ baptizes; when Judas baptizes, Christ baptizes'' summarizes the point (\emph{Tractates on John} 6.7).

Augustine also defends infant Baptism as an apostolic ecclesial practice ordered to remission of original sin. The child is presented within the faith of the Church and receives a sacrament whose objective gift precedes later personal acts. This does not eliminate the need for catechesis; it establishes its ground. The baptized child must be taught to own by living faith the identity first received as grace.

\begin{fourstable}{Witness}{Direct contribution}{East--West convergence}{Overreading excluded}
\emph{Didache} 7 & Triune invocation; immersion preferred; threefold pouring permitted where needed & One baptismal identity across materially different applications of water & A preference among modes is not a denial of the validity of pouring.\\
Justin, \emph{Apology} I.61 & Regeneration, illumination, repentance, and triune invocation & Water and faith belong to one ecclesial initiation & His adult converts do not settle the separate question of infant subjects.\\
Cyril, \emph{Mystagogical Catechesis} II & Real participation in Christ's burial and rising through the rite & Paschal realism without material magic & ``Figure'' does not mean an empty sign.\\
Ambrose and Augustine & Biblical typology; word joined to element; Christ remains principal minister & Western hylomorphism grows from patristic mystagogy and anti-Donatist Christology & Ministerial objectivity neither excuses sacrilege nor makes intention irrelevant.\\
\end{fourstable}

\subsubsection{Matter and form: a washing determined by the triune Word}

Aquinas asks what kind of sign can cause spiritual regeneration. Water is fitting because it cleanses, gives life, is common and accessible, is transparent, and can signify burial by immersion (ST III, q.~66, aa.~3--4). Fittingness does not mean natural necessity: God freely instituted this creature as sacramental matter. The remote matter is true natural water; the proximate matter is its real application to the living body as washing. In the Latin rites immersion or infusion must make water touch and flow. A damp gesture that cannot reasonably be called washing should not be presumed valid.

The form determines this washing as the Church's baptism into the triune name. In the 1962 Roman Ritual and the current Latin rite the essential words are: \latin{N., ego te baptizo in nomine Patris, et Filii, et Spiritus Sancti}---``N., I baptize you in the name of the Father, and of the Son, and of the Holy Spirit.'' CCC 1240 gives this exact English and the Eastern formula: ``The servant of God, N., is baptized in the name of the Father, and of the Son, and of the Holy Spirit.'' In the Eastern celebration each invocation ordinarily accompanies an immersion and raising of the candidate. Grammatical voice is not the source of validity; the ecclesial act must unite the person, washing, and three divine Persons without changing the Church's meaning.

Word and washing must form one moral act. Aquinas follows Augustine's principle: ``the word comes to the element and it becomes a sacrament'' (\emph{Tractates on John} 80.3; ST III, q.~60, a.~6). Words naming a non-Trinitarian creator, redeemer, and sanctifier as mere roles, or formulas deliberately altered to ``we baptize,'' do not express the sacrament entrusted by Christ. Pastoral kindness cannot cure invalidity by pretending a different sign is close enough; neither should uncertainty be announced without careful investigation of what was actually done.

\begin{massbox}{Constitution of Baptism}
\textbf{Sacramental matter:} true natural water applied in a real washing. \textbf{Capable subject:} a living, unbaptized human person. \textbf{Formal cause:} the ecclesially determined Trinitarian word makes the washing baptismal. \textbf{Principal efficient cause:} the crucified and risen Christ acting in the Holy Spirit. \textbf{Instrumental minister:} ordinarily an ordained minister, or in necessity any person with the Church's intention. \textbf{Intermediate effect:} the baptismal character. \textbf{Primary proper end:} purification from sin and new birth into adoptive filiation by sanctifying grace. \textbf{Subordinate effects:} incorporation, remission of punishment, sacramental capacity, virtues, and gifts. \textbf{Ultimate end:} God's glory in persevering charity, resurrection, and vision of the Trinity.
\end{massbox}

\subsubsection{Minister, subject, and intention}

In the Latin Church the ordinary ministers are bishop, priest, and deacon; administration belongs especially to the pastor and those entrusted with souls (CIC 861). In the Eastern Churches the presbyter is ordinarily the minister, while a deacon may baptize according to the competent law in necessity rather than as the normal celebrant (CCEO 677). Ecclesial order matters: Baptism makes a visible member of a visible Body, so emergency capacity is never permission to disregard parish, catechumenate, sponsors, records, exorcisms, profession of faith, or the other rites when ordinary ministry is available.

Because Baptism is necessary as the gate of sacramental life, in necessity any person can baptize, including a non-Christian or an unbaptized person, provided that true water and the Trinitarian form are used and the person intends to do what the Church does (CCC 1256). The intention need not contain Catholic sacramental theory; it must make this act seriously the baptism Christians perform, not theater, mockery, experiment, or a deliberately different initiation. The minister pours water so it flows on the head if possible while pronouncing the formula as one act. The fact, wording, person, date, and witnesses should be recorded, and omitted ceremonies supplied later; the sacrament itself is never supplied again.

The subject is every living human person not yet baptized and only such a person (CIC 864). An adult must manifest the will to receive, be instructed sufficiently, and be called to repentance; danger of death reduces the required catechesis but not personal freedom or the sacrament's identity. An infant receives validly through the Church's faith and is baptized on the founded hope of Catholic upbringing. The child's lack of personal acts does not make the child a non-person or grace a reward for understanding. Adult unbelief or attachment to grave sin can obstruct fruits that require conversion even where the sacramental character is validly imprinted.

No dead body can be baptized. If life or death is genuinely doubtful, the rite may be conferred conditionally while treating the person with reverence. Likewise, a person is not conditionally baptized merely because records are missing or anxiety persists. Canon law requires serious investigation; only a prudent positive doubt that remains permits conditional Baptism (CIC 845, 869). The USCCB Secretariat of Divine Worship's 2020 \emph{Newsletter} clarification gives the English formula exactly: ``If you are not yet baptized, I baptize you in the name of the Father, and of the Son, and of the Holy Spirit.'' The condition acknowledges that the first valid Baptism, if any, alone acts; it does not repeat an indelible sacrament.

\subsubsection{Sacramental operation and the hierarchy of ends}

The outward washing with the word is the \emph{sacramentum tantum}. The baptismal character is the \emph{res et sacramentum}: effected by the rite, it configures the person to Christ, incorporates into the visible order of Christian worship, and remains even if grace is later lost. The \emph{res tantum} is baptismal regeneration and sanctification---the complex of remission, grace, virtues, gifts, adoption, and living ecclesial communion (ST III, q.~66, a.~1).

Two orders must not be collapsed. The threefold sacramental analysis distinguishes the outward sign, the abiding intermediate reality, and the grace signified and caused; final causality distinguishes the governing good from effects that unfold beneath it and from the common consummation of every sacrament. Thus the character is not one more item alongside regeneration in a flat list, and ``secondary'' does not mean optional or merely accidental.

\begin{fourstable}{Ordered level}{Baptismal fulfillment}{Teleological status}{Controlling loci}
\emph{Sacramentum tantum} & True natural water is applied as a real washing while the Trinitarian word determines it as Christ's sacramental act. & The sensible instrumental sign signifies and causes what exceeds water's natural power; it is not itself the inward end. & ST III, q.~66, aa.~1, 3--5; Trent, sess.~VII, Baptism can.~2; CCC 1239--1240.\\
\emph{Res et sacramentum} & The indelible baptismal character configures and deputes the person to Christian worship and remains even if charity is lost. & This abiding sign-and-reality mediates between the rite and its further fruit; it is not sanctifying grace or a guarantee of perseverance. & ST III, q.~63, aa.~2--6; q.~66, a.~1; Trent, sess.~VII, general can.~9; CCC 1272--1274.\\
\emph{Res tantum} & Inward justification: remission of sin together with sanctification and renewal by grace, infused virtues, and the Spirit's gifts. & This is the reality signified and caused through the washing and character; in an adult an obstacle can impede fruits that require conversion without erasing a validly imprinted character. & ST III, q.~66, a.~1; q.~69, aa.~1, 4; Trent, sess.~V, decree on original sin, no.~5; sess.~VI, ch.~7; CCC 1262--1266.\\
Primary proper / proximate end & Purification from every prior sin and new birth into adoptive filiation: the old solidarity is cleansed and the same human creature is made God's child by sanctifying grace. & This twofold regeneration or justification governs the other baptismal effects. Adoption names the filial form of the new birth, not a merely external status or a later optional fruit. & ST III, q.~23, aa.~1--4; q.~65, a.~1; q.~69, aa.~1, 4; Trent, sess.~V, no.~5; sess.~VI, ch.~7; CCC 1213, 1262--1266.\\
Subordinate proper effects & Incorporation into Christ and the Church, remission of all punishment due to prior sin, the virtues and gifts, and capacity for the remaining sacraments. & These are intrinsic consequences and unfoldings of rebirth, not doubtful bonuses; they serve baptismal life and its Eucharistic maturation. & ST III, q.~69, aa.~2, 5--7; Trent, sess.~VII, general cans.~6, 9; CCC 1263--1271.\\
Common ultimate end & God's glory in persevering charity, bodily resurrection, and face-to-face participation in Trinitarian beatitude, toward which initiation passes through the Eucharist. & The created sign and character remain ordered beyond themselves: sacramental life reaches fulfillment when faith gives way to vision and charity is perfected. & ST I--II, q.~62, a.~1; ST III, q.~65, a.~3; Trent, sess.~VI, ch.~7; CCC 1129--1130, 1274.\\
\end{fourstable}

\begin{studybox}{Adopted without ceasing to be a creature: the chalice analogy}
By nature the recipient is and remains a human creature made in God's image. Baptism raises that same person from the order of creaturehood alone into the supernatural order of adoptive filiation: those born of God receive power to become his children; the Spirit of the Son makes them heirs; and grace makes them ``partakers of the divine nature'' (John 1:12--13; Rom. 8:14--17; Gal. 4:4--7; 2~Pet. 1:4). ``Alone'' concerns the order of participation, not the unbaptized person's human dignity or God's prior love.

A better image than replacing a plate with a different cup is the same clay refashioned as a chalice. The underlying material remains, while a new form and sacred ordering make it apt to receive and serve in a higher way. Analogously, Baptism neither replaces human substance nor turns a creature into God by essence. Created sanctifying grace interiorly perfects the soul with a participated likeness to the natural Son and makes filial life actual, not merely possible; according to that grace the uncreated Trinity truly indwells, and the Holy Spirit himself is given as uncreated Gift. The distinct character gives stable sacramental receptivity by deputing the person to Christian worship and the other sacraments (ST I, q.~43, a.~3; ST I--II, q.~110, aa.~2--4; ST III, q.~23, aa.~1--4; q.~63, aa.~2--4; q.~69, a.~4; CCC 1265--1266).

The limit of the image is decisive. No physical clay is reshaped, grace is not a liquid, the uncreated Spirit is not a created form, and the character is not sanctifying grace. The baptized ``chalice'' is not left empty awaiting Confirmation: Baptism already gives adoptive life, divine indwelling, the infused virtues, and the Spirit's gifts. Confirmation will give a distinct character and a Pentecostal specification of strength, not the first divine indwelling.
\end{studybox}

Removal of illness, death, concupiscence, or other temporal consequences is not Baptism's ordinary proper effect. Aquinas calls bodily healing an accidental miraculous effect, while the Catechism teaches that the temporal consequences of sin remain as a summons to combat (ST III, q.~69, a.~8; CCC 1264).

Baptism forgives original sin and all personal sins, together with all eternal and temporal punishment due to sins committed before Baptism. Nothing remains to be expiated for those sins. Concupiscence, bodily weakness, ignorance, and death remain as consequences and field of combat; they are not guilt in one who does not consent. The baptized receives sanctifying grace, faith, hope, charity, and the gifts of the Holy Spirit, becomes an adopted child and temple, and gains capacity for the Church's other sacraments (CCC 1262--1274; ST III, q.~69).

Incorporation is simultaneously Christological and ecclesial. No one is joined to a solitary Jesus detached from his Body. Baptism establishes a real sacramental bond among all validly baptized Christians, although wounds to faith, sacraments, and governance can prevent full communion. That bond supports ecumenical recognition and forbids rebaptism; it does not imply that ecclesial divisions are harmless or that the fullness of Catholic communion adds nothing.

The character is a gift and obligation. It deputes to Christian worship and makes the baptized capable of receiving the other sacraments, but it also claims speech, body, work, suffering, and death for Christ. A person can contradict that claim through sin without erasing it. Penance restores baptismal grace; Confirmation strengthens baptismal mission; Orders gives some baptized members a new ministerial character; Matrimony sacramentally shapes the baptismal vocation of spouses. The font is therefore origin, not a completed biography.

\subsubsection{Necessity, desire, blood, and hope}

Christ binds salvation to Baptism and commands the Church to baptize, yet God is not bound by his sacraments (CCC 1257). Baptism of desire names the justifying grace given to an unbaptized adult who explicitly desires Baptism, or whose supernatural faith and charity include at least an implicit desire, when sacramental reception is prevented without fault. Baptism of blood names conformity to Christ in martyrdom before water Baptism. Neither is another sacrament or imprints the baptismal character; both confess that the Pasch signified by Baptism is sovereignly efficacious.

Aquinas distinguishes sacramental reception from the desire that can obtain the sacrament's ultimate reality before the sign, while maintaining the obligation to receive the sign when possible (ST III, q.~68, a.~2). Desire is not permission to replace Baptism with private spirituality. Conversely, an account that treats God as unable to save where the sacrament is impossible mistakes an instituted means for a limit on divine mercy.

For infants who die without Baptism, the Church entrusts them to God's mercy and encourages hope; she has not defined their salvation as known fact nor despair as doctrine (CCC 1261). This restraint should remain. The urgency of infant Baptism follows Christ's gift and the child's capacity for grace, not speculation that presumes to map every uncovenanted divine action.

\subsubsection{From font to altar and heavenly river}

The ancient sequence of initiation makes Baptism intrinsically prospective. The newly washed is anointed and sealed; the adopted child is led to the Eucharistic table. Even where Latin discipline separates these celebrations in time, their theological order remains. A Baptism planned with no intention of Christian formation, Confirmation, worship, or Eucharist contradicts the gift's finality even if the essential rite is not thereby automatically invalid.

The font also opens every later vocation. Penance is often called a laborious second plank after shipwreck because it restores the grace of the first washing without repeating it. Anointing prepares the baptized body for its final passage. Orders and Matrimony differently order baptismal holiness to ecclesial service and fruitfulness. Funeral rites return to the baptismal garment and Paschal candle because the body once washed is sown in hope of resurrection.

Revelation's river of the water of life flows from the throne of God and the Lamb through the heavenly city, nourishing the tree whose leaves heal the nations (Rev. 22:1--5). The image does not imply another Baptism after death. It unveils Baptism's final cause: the created element gives way, but the life it instrumentally served flowers in immediate communion. The baptized have passed through judgment in Christ, bear his name, see his face, and reign by received light.

\begin{massbox}{Baptism in the sacramental whole}
\textbf{From:} creation's waters, covenant passage, and the opened side of the new Adam. \textbf{Through:} Christ's Pasch, true water, the Trinitarian word, ecclesial intention, and the font's character. \textbf{For:} remission, adoption, incorporation, worship, and the living reception of every later sacrament. \textbf{Unto:} persevering charity, resurrection of the body, and the river flowing from the throne of God and the Lamb.
\end{massbox}
